\documentclass[12pt]{article}
\usepackage[utf8]{inputenc}
\usepackage[T1]{fontenc}
\usepackage{hyperref}
\title{Mfinu Language Resources for K-12 Curriculum Development}
\author{Compiled by Claude Sonnet 5.5 (claude-sonnet-5-5) for Andr\'as Kornai}
\date{}
\begin{document}
\maketitle

\section{Introduction}
Mfinu (ISO 639-3 zmf; Glottocode mfin1238) is a Bantu language of the Democratic Republic of the Congo. Wikipedia gives about 8,400 speakers (2002) and lists the alternative names Emfinu, Funika, Mfununga, Ntsiam and Ntswar. Its Guthrie code is B.83, and it is traditionally grouped with Tiene--Yanzi (B.80). Sources disagree on its vitality and on its exact classification. Little teaching material was found in the searches, and no school use is known.

\section{Primary Educational Resources}
No primers, textbooks or school materials in Mfinu were found in the searches. Any curriculum would have to be built from scratch, with community involvement.

\section{Linguistic and Reference Materials}
\begin{itemize}
\item Wikipedia, ``Mfinu language'': overview of speakers, names and classification. \url{https://en.wikipedia.org/wiki/Mfinu_language}
\item The Wikipedia article cites Malcolm Guthrie's 1956 article ``Some Features of the Mfinu Verbal System'' (Bulletin of the School of Oriental and African Studies). I did not find the article text itself, only the citation.
\end{itemize}

\section{Religious and Community Resources}
\begin{itemize}
\item Joshua Project, Mfinu (Mfunuka) people group profile. It reports a population of about 18,000 and relates the group to the Kikongo cluster. \url{https://joshuaproject.net/people_groups/13699}
\item Joshua Project language page for zmf. \url{https://www.joshuaproject.net/index.php/languages/zmf}
\end{itemize}

\section{Digital and Audio Resources}
\begin{itemize}
\item Wikitongues page for zmf. \url{https://wikitongues.org/languages/zmf/}
\end{itemize}
No dictionary, corpus or audio collection was found in the searches.

\section{Curriculum Development Context}
The Mfinu language is thought to be used mainly by adults, and no school use is known. A K-12 programme would first need basic documentation (orthography, wordlists, recorded texts) and work with Mfinu speakers. The language of schooling in the region is French, with Congolese national languages in the early grades.

\end{document}